\documentclass[11pt,a4paper]{article}
\usepackage{babel}
\babelprovide[import,main]{spanish}
\usepackage{fontspec,amsmath,amssymb,booktabs,geometry}
\geometry{margin=2.2cm}
\setmainfont{Latin Modern Roman}
\setlength{\parindent}{0pt}
\setlength{\parskip}{6pt}
\title{Ejercicio 1 del set\\Desarrollo para comparar a mano y en PRE-POST 2.0}
\author{Guía de cálculo revisada — no acredita trabajo manuscrito}
\date{Entrega 2 · 27 de octubre de 2026}
\begin{document}
\maketitle
\section{Enunciado, alcance y adopciones}
Fuente: \textit{SET DE EJERCICIOS PARA AVANCE 2.pdf}, ejercicio 1, páginas 1--3 (UNaM, Dávalos y Fraticelli). Se analiza una sección rectangular en flexión simple. Los datos del enunciado son:
\[
b=0.40\,\mathrm m,\quad h=0.90\,\mathrm m,\quad d_p=0.80\,\mathrm m,
\]
\[
f'_c=30\,\mathrm{MPa},\quad M_s=1000\,\mathrm{kN\,m},\quad f_{pe}=1150\,\mathrm{MPa},\quad f_{py}=1700\,\mathrm{MPa}.
\]
La fuente utiliza acero grado 270 y $M_u=1.6M_s$. Para el procedimiento resistente implementado se adopta explícitamente $f_{pu}=1860\,\mathrm{MPa}$. Se adopta la fuerza efectiva por exceso en incrementos de $10\,\mathrm{kN}$: es una decisión de esta resolución, no un dato del enunciado.

No se supone luz, carga distribuida ni pérdidas de 15\%. $f_{pe}$ ya es la tensión efectiva del acero. Este ejercicio no verifica transferencia, pérdidas, corte ni flechas; esas verificaciones quedan fuera del alcance actual de la app.

\section{Propiedades y convención de signos}
Compresión negativa, tracción positiva; excentricidad positiva hacia arriba y momento sagante positivo.
\begin{align*}
A&=bh=0.40(0.90)=0.36\,\mathrm{m^2},\\
I&=\frac{bh^3}{12}=\frac{0.40(0.90)^3}{12}=0.0243\,\mathrm{m^4},\\
W&=\frac{I}{h/2}=\frac{0.0243}{0.45}=0.054\,\mathrm{m^3},\\
e&=\frac h2-d_p=0.45-0.80=-0.35\,\mathrm m.
\end{align*}
Las tensiones elásticas extremas son
\[
\sigma_{sup}=-\frac PA-\frac{Pe+M_s}{W},\qquad
\sigma_{inf}=-\frac PA+\frac{Pe+M_s}{W}.
\]
\newpage
\section{Fuerza efectiva y área activa}
El intervalo de clase T utilizado en la fuente, con $f'_c$ en MPa, es abierto:
\[
0.7\sqrt{30}<\sigma_{inf}<\sqrt{30},\quad
3.834057903<\sigma_{inf}<5.477225575\quad[\mathrm{MPa}].
\]
Trabajando en N, m y Pa:
\[
\sigma_{inf}=\frac{10^6}{0.054}-P\left(\frac1{0.36}+\frac{0.35}{0.054}\right).
\]
El coeficiente de $P$ vale $9.259259259\,\mathrm{m^{-2}}$. Los límites de fuerza son:
\begin{align*}
P_{inf}&=\frac{10^6/0.054-\sqrt{30}\,10^6}{9.259259259}
=1408.459637894\,\mathrm{kN},\\
P_{sup}&=\frac{10^6/0.054-0.7\sqrt{30}\,10^6}{9.259259259}
=1585.921746526\,\mathrm{kN}.
\end{align*}
Por tanto, $1408.459638<P<1585.921747$ kN. La igualdad en el límite no se clasifica como T. Adoptando por exceso al siguiente múltiplo de 10 kN:
\[
\boxed{P=1410\,\mathrm{kN}},\qquad
A_p=\frac P{f_{pe}}=\frac{1410000}{1150}=\boxed{1226.086956522\,\mathrm{mm^2}}.
\]
Esta área es teórica. No se selecciona una cantidad de torones sin definir el área por torón y su distribución.
\section{Tensiones de servicio}
\begin{align*}
\frac PA&=\frac{1410000}{0.36}=3.916666667\,\mathrm{MPa},\\
Pe+M_s&=1410000(-0.35)+1000000=506500\,\mathrm{N\,m},\\
\frac{Pe+M_s}{W}&=\frac{506500}{0.054}=9.379629630\,\mathrm{MPa},\\
\sigma_{sup}&=-3.916666667-9.379629630=\boxed{-13.296296296\,\mathrm{MPa}},\\
\sigma_{inf}&=-3.916666667+9.379629630=\boxed{5.462962963\,\mathrm{MPa}}.
\end{align*}
Se satisface $3.834057903<5.462962963<5.477225575$: la fibra inferior queda en clase T según el criterio de la fuente. Esto no significa que todas las verificaciones resistentes estén satisfechas.
\newpage
\section{Resistencia a flexión: procedimiento revisado}
El software utiliza el procedimiento ya implementado de CIRSOC 201-2005, 18.7.2(a) y 9.3.2.1. Es distinto de la simplificación del desarrollo impreso. Su adopción académica y la unificación con el material ACI del curso siguen pendientes de ratificación docente.

Se trabaja en mm, N y MPa: $b=400$, $d_p=800$, $A_p=1226.086956522$.
\begin{align*}
M_u&=1.6M_s=\boxed{1600\,\mathrm{kN\,m}},\\
\beta_1&=0.85,\qquad \gamma_p=0.28\quad(f_{py}/f_{pu}=1700/1860\geq0.90),\\
\rho_p&=\frac{A_p}{bd_p}=\frac{1226.086956522}{400(800)}=0.003831521739,\\
f_{ps}&=f_{pu}\left[1-\frac{\gamma_p}{\beta_1}\frac{\rho_pf_{pu}}{f'_c}\right]\\
&=1860\left[1-\frac{0.28}{0.85}\frac{0.003831521739(1860)}{30}\right]
=1714.449053708\,\mathrm{MPa},\\
T&=A_pf_{ps}=2102063.622373\,\mathrm N,\\
a&=\frac{T}{0.85f'_cb}=\frac{2102063.622373}{0.85(30)(400)}=206.084668860\,\mathrm{mm},\\
c&=\frac a{\beta_1}=242.452551600\,\mathrm{mm},\\
\varepsilon_t&=0.003\frac{d_p-c}{c}=0.006898844059>0.005.
\end{align*}
El criterio implementado admite $\phi=0.90$ en este dominio. El brazo interno es
\[
z=d_p-a/2=800-206.084668860/2=696.957665570\,\mathrm{mm}.
\]
\begin{align*}
M_n&=Tz=1465049355.129\,\mathrm{N\,mm}=1465.049355129\,\mathrm{kN\,m},\\
\phi M_n&=0.90(1465.049355129)=\boxed{1318.544419616\,\mathrm{kN\,m}}.
\end{align*}
\[
\boxed{\phi M_n=1318.544420<1600=M_u:\quad\text{NO CUMPLE}}
\]
El objetivo de la comparación es reproducir correctamente este resultado, no modificar los datos hasta obtener cumplimiento. Si se pide rediseñar, será una etapa separada.
\newpage
\section{Por qué no copiamos los números impresos}
La guía publica $P=1403$ kN y $A_p=1220$ mm$^2$. Sin embargo:
\[
\sigma_{inf}(1403)=\frac{1000000}{0.054}-1403000\left(\frac1{0.36}+\frac{0.35}{0.054}\right)
=5.527777778\,\mathrm{MPa}>\sqrt{30}.
\]
La fuerza impresa no satisface el límite superior de clase T con los datos publicados. Además, la resolución resistente de la guía usa una simplificación diferente. No debe interpretarse la diferencia de resistencia como error de redondeo del programa.

\begin{center}
\begin{tabular}{lrr}
\toprule
Magnitud & Guía impresa & Resolución revisada\\
\midrule
$P$ (kN)&1403&1410\\
$A_p$ (mm$^2$)&1220&1226.086957\\
$M_u$ (kN m)&1600&1600\\
\bottomrule
\end{tabular}
\end{center}
\section{Cómo ejecutar y comparar}
\begin{enumerate}
\item Iniciar la app con el comando:\\ \texttt{python -m streamlit run app/main.py}.
\item La app abre directamente el ejercicio 1 del set, su único caso.
\item Pulsar ``Restablecer ejercicio del set'' y ``Calcular ejercicio 1''.
\item Desarrollar a mano las secciones 2--5 de este documento, conservando cifras intermedias.
\item Ingresar los resultados propios en ``Comparar con tu desarrollo manuscrito'', respetando las unidades; pulsar ``Comparar resultados''.
\item Descargar el cotejo JSON y la memoria PDF. Adjuntar el manuscrito real y registrar quién lo elaboró y revisó.
\end{enumerate}
El error absoluto es $|v_s-v_m|$; el relativo es $100|v_s-v_m|/|v_m|$, donde $v_s$ es el valor del software y $v_m$ el manual. Para referencia cero se informa solo el absoluto. La app no acredita autoría ni revisión humana automáticamente.

\begin{center}
\begin{tabular}{lrrl}
\toprule
Magnitud & Referencia revisada & Mi cálculo & Unidad\\
\midrule
$P$&1410&\rule{1.5cm}{0.2pt}&kN\\
$A_p$&1226.086957&\rule{1.5cm}{0.2pt}&mm$^2$\\
$\sigma_{sup}$&$-13.296296$&\rule{1.5cm}{0.2pt}&MPa\\
$\sigma_{inf}$&5.462963&\rule{1.5cm}{0.2pt}&MPa\\
$M_u$&1600&\rule{1.5cm}{0.2pt}&kN m\\
$\phi M_n$&1318.544420&\rule{1.5cm}{0.2pt}&kN m\\
\bottomrule
\end{tabular}
\end{center}
\end{document}
